\begin{table}[ht]
\centering
\caption{Dynamic-design sensitivity summary.}
\label{tab:dynamic_sensitivity_summary}
\begin{threeparttable}
\small
\setlength{\tabcolsep}{5pt}
\begin{tabular}{p{.42\linewidth}rp{.34\linewidth}}
\toprule
Diagnostic quantity & Value & Interpretation \\
\midrule
BJS coefficient at first post period & 0.052 & First post-exposure timing estimate \\
Standard error & 0.018 & BJS period-specific standard error \\
Maximum absolute pre-period coefficient & 0.041 & Observed pre-period deviation scale \\
Breakdown linear-extrapolation $M$ & 0.018 & Smallest linear-extrapolation allowance that reaches zero \\
Survives observed pre-period maximum & no & Diagnostic robustness indicator \\
\bottomrule
\end{tabular}
\begin{tablenotes}[flushleft]\footnotesize
\item \textit{Notes:} The BJS event study is used to assess timing patterns, not as the primary identifying design. The final row records whether the first post-period estimate remains different from zero when the allowed post-period violation equals the observed maximum absolute pre-period coefficient. The sensitivity calculation uses the saved BJS period-specific standard errors; the off-diagonal covariance is unavailable in the source BJS event-study output.
\end{tablenotes}
\end{threeparttable}
\end{table}

